\section*{Nomenclature}
\addcontentsline{toc}{section}{Nomenclature}

\noindent\textbf{Frames and mechanics (\cref{sec:mechanical,sec:simulation})}
\begin{center}
\begin{tabular}{@{}p{0.2\linewidth}p{0.74\linewidth}@{}}
$x$, $y$ & robot frame with origin at the midpoint of the axle: $x$ forward, $y$ to the left \\
$m$, $W$ & mass and weight of the robot \\
$x_c$, $x_k$ & distance of the centre of mass and of the caster ahead of the axle \\
$N_w$, $N_k$ & static load on the two wheels together and on the caster \\
$\mu$ & coefficient of friction \\
$a_{\max}$ & largest acceleration that the wheels can transmit without slip \\
$d$, $h$ & distance of the centre of mass from an edge of the support triangle, its height above the floor \\
$r$, $b$ & wheel radius and wheel separation \\
$\omega_l$, $\omega_r$ & angular speed of the left and the right wheel \\
$v$, $\omega$ & forward speed and turn rate of the robot \\
$f$ & focal length of the parabolic jaw profile \\
$\alpha$ & angle between the jaw surface and the direction of travel \\
$\varphi$ & friction angle, $\arctan\mu$ \\
$h_c$, $x_{\min}$ & height of the camera above the floor, nearest floor point that it sees \\
\end{tabular}
\end{center}

\noindent\textbf{Line following and mission (\cref{sec:arena,sec:autonomy,sec:results})}
\begin{center}
\begin{tabular}{@{}p{0.2\linewidth}p{0.74\linewidth}@{}}
$p$ & pitch of the infrared channels \\
$V_i$, $V_{\mathrm{ref}}$ & voltage of channel $i$, voltage on a white floor \\
$d_i$ & darkness of channel $i$, from 0 on white to 1 on black \\
$e$ & offset of the line from the centre of the array, positive to the left \\
$k_p$, $k_d$ & proportional and derivative gain of the steering law \\
$\omega_{\max}$ & limit of the commanded turn rate \\
$L$ & distance of the array ahead of the axle \\
$\psi$ & heading of the robot relative to the line \\
$\omega_n$, $\zeta$ & natural frequency and damping ratio of the line-following loop \\
station pose & pose in front of a stem at which the navigation hands over to the line follower \\
marker & floor marking wider than the array: the start square, a crossbar or a drop square \\
\end{tabular}
\end{center}

\noindent\textbf{Abbreviations}
\begin{center}
\begin{tabular}{@{}p{0.2\linewidth}p{0.74\linewidth}@{}}
CAD & computer-aided design \\
EKF & extended Kalman filter \\
IMU & inertial measurement unit \\
IR & infrared \\
PD & proportional and derivative (control) \\
RMS & root mean square \\
SLAM & simultaneous localisation and mapping \\
URDF & Unified Robot Description Format \\
\end{tabular}
\end{center}
\clearpage
